\scanpage{74}
% archive sheet number: 73

\handwritten{7.} Nous verrons çi-dessus que la réponse à la question précédente est affirmative lorsque $k$ est un corps fini, en particulier si c'est le \unsure{corps|coops} premier $\underline{F}_p$, On trouve alors, grâce à cette structure d'extension sur $\underline{G}(\underline{F}_p)$, un homomorphisme canonique de Gerbes (pour tout $k$)
\[
\underline{G}(k)\handwritten{\simeq}\underline{G}(\underline{F}_p)\mathrm{x}_{\deleted{\unsure{\underline{G}(}}\handwritten{\unsure{\underline{\pi}|\underline{\Pi}}}_p(\underline{F}_p))}\handwritten{\underline{\pi}}_p(k)\ ,
\]
et pour $k$ parfait, une réponse affirmative à la question précédente revient à l'assertion que l'homomorphisme précédent est une équivalence de Gerbes. Comme on connait parfaitement bien la Gerbe $\underline{G}(\underline{F}_p)$ et son homomorphisme dans $\handwritten{\underline{\pi}}_p(\underline{F}_p)$ (cf plus bas) cela donnerait donc la structure de $\underline{G}(k)$ en termes de la connaiassance de \deleted{\unsure{$\pi_p(\underline{F}_p$}} $\handwritten{\underline{\pi}}_{\deleted{\unsure{a}}p}(k)$ (connu quand\handwritten{\,}on connait $\handwritten{\pi}_1(k)=\operatorname{Gal}(\handwritten{\bar{k}}/k)$) et \deleted{de} l'homomorphisme $\handwritten{\underline{\pi}}_p(k\deleted{\unsure{'}})\handwritten{\longrightarrow\underline{\pi}}_p(\underline{F}_p)$ (connu quand on connait $\handwritten{\pi}_1(k)\handwritten{\longrightarrow\pi}_1(\underline{F}_p)\handwritten{\longrightarrow}\handwritten{\hat{\underline{Z}}}$). Cela résoudrait donc la question de la structure de $\underline{G}(k)$ pour $k$ parfait. On notera que l'objet de Tate de $\mathrm{Fc}\handwritten{\mathrm{r}}\mathrm{iso}(k)$ provient de celui de $\mathrm{Fcriso}(\underline{F}_p)$ et sera donc également déterminé. Le foncteur fibre canonique sur $K$\inserted{$=K(k)$} correspond alors à un foncteur fibre de $\mathrm{Fcriso}(\underline{F}_p)$ sur $K$, un foncteur fibre de \deleted{\unsure{$\underline{\pi}_p(k)$}}\inserted{la catégorie des $\underline{Q}_p$-faisceaux constr. sur $k$} sur $K$, et un isomorphisme entre les ``restrictions'' de ces foncteurs à la catégorie des $\underline{Q}_p$-faisceaux constructibles sur $\underline{F}_p$ ; or le premier foncteur fibre n'est autre que \inserted{celui déduit par} l'extension des scalaires $\underline{Q}_p=K\handwritten{(}\underline{F}_p)\handwritten{\longrightarrow}K=K(k)$ du foncteur fibre analogue sur $\unsure{\mathbb{F}|\mathrm{F}|\mathrm{G}}\mathrm{criso}(\underline{F}_p)$, le deuxième est \unsure{le|he} foncteur \inserted{\handwritten{$i$:}} $V\handwritten{\mapsto}V\otimes_{\underline{Q}_p}K$, et l'isomorphisme \deleted{canonique} entre les restrictions est l'isomorphisme canonique évident\inserted{\handwritten{(compatibilité de $i$ aux changements de corps de base)}}. Donc la structure supplémentaire de $\underline{G}(k)$ (savoir sa section \deleted{sur} canonique sur $K$) est déterminée par la connaissance de la section de $\handwritten{\underline{\pi}}_p(k)$ sur $K$ (i.e. le foncteur précédent $V\handwritten{\mapsto}V\otimes_{\underline{Q}_p}K$ ) et l'isomorphisme de l'\unsure{image|iamge} de ladite section dans $\handwritten{\underline{\pi}}_p(\underline{F}_p)$ avec la section provenant de\deleted{\unsure{s}} la section canonique de $\handwritten{\underline{\pi}}_p(\underline{F}_p)$ sur $\underline{Q}_p$ ; ce sont là des données qui ne sont évidemment pas contenues dans celle du groupe

\margin{\handwritten{\uline{NB} On a donc, en comparant le foncteur ``fibre \deleted{``\unsure{oubli}} en $\bar{k}$'' sur $Q_p$, et le foncteur précédent, un torseur canonique sur $K$ de groupe \deleted{\unsure{de}}\unsure{lien} $\underline{\pi}_p(k)_K$\deleted{\illegible{2} $\unsure{Q_p|\mathbb{Q}_p}$}}}
